  % !TeX root = book.tex
\chapter{Prolog AI agents on the Agentic Web}
\label{sec:prolog-ai-agents-on-the-agentic-web}

The chart in Figure~\ref{fig:ai-and-prolog} describes the evolution of the field of Artificial Intelligence since the late 1950s and up till now. The evolution is described as cyclic, where AI ``booms'' (or ``summers'') are followed by AI ``busts'' (or ``winters''). According to this view we have summer now and the excitement over AI has peaked once again and to a height never before seen.

\begin{figure}[h]
    \centering
	\includegraphics[width=11cm]{ai_booms.jpg}
    \caption{Cycle of AI booms and busts. Chart inspired by Yutaka Matsuo.}
    \label{fig:ai-and-prolog}
\end{figure}

\noindent Historically, Prolog and Lisp played prominent roles during the second boom, where the \textit{expert system} was seen as the flagship technology of AI application. As we entered the third boom, the focus shifted toward non-symbolic methods, such as deep learning and Large Language Models (LLMs). However, this chapter argues that the story does not end there. We draw on the rapidly spreading notion of the \textit{Agentic Web}, an evolution of the World Wide Web into an environment where software agents operate alongside traditional resources. While the classic Web serves primarily as a distributed hypertext system for human consumption, the Agentic Web envisions a dynamic society of autonomous software processes capable of perceiving, reasoning, and acting within the web space.

This vision intersects naturally with logic programming. Prolog's declarative foundations -- combined with its suitability for knowledge representation, inference, and dynamic communication -- make it an apt candidate for implementing agents that are both knowledgeable and capable of structured interaction. We call this the \textit{Agentic Prolog Web}.

This chapter unfolds in four progressive phases. Section~\ref{sec:symbolic-prolog} establishes the \textit{symbolic competence} of Prolog by revisiting classic AI idioms on a single machine. Section~\ref{sec:prolog-web-agents} \textit{distributes} those idioms across the Agentic Web, transforming predicates into message-driven services. Section~\ref{sec:embodied-agents} explores \textit{user-interface} (UI) agents and embodied conversational agents. Finally, Section~\ref{sec:neuro-symbolic} extends this substrate with neural components, showing how Web Prolog can orchestrate Large Language Models to form robust \textit{neuro-symbolic} agents.

===

Figure~\ref{fig:ai-and-prolog} sketches a popular retrospective view of Artificial Intelligence as a sequence of booms and busts: periods of intense optimism and investment followed by ``AI winters''. On this view, we are currently in a new boom whose scale and public visibility surpass earlier waves.

\begin{figure}[h]
\centering
\includegraphics[width=11cm]{ai_booms.jpg}
\caption{Cycle of AI booms and busts. Chart inspired by Yutaka Matsuo.}
\label{fig:ai-and-prolog}
\end{figure}

%\begin{figure}[h]
%\centering
%\includegraphics[width=\textwidth]{ai-timeline.jpg}
%\caption{Cycle of AI booms and busts. }
%\label{fig:ai-and-prolog}
%\end{figure}

\noindent Prolog and Lisp were prominent during earlier waves, especially the expert-systems era, whereas the current boom is driven primarily by data-intensive methods such as deep learning and large language models. This chapter argues that the story does not end there. As software becomes increasingly \emph{agentic}, the need for controllable behaviour and auditable reasoning re-emerges as a first-class engineering concern.

Earlier in the book, and in particular in Chapter~\ref{sec:the-prolog-web}, we argued that the Prolog Web already constitutes a multi-agent system in a minimal but precise sense: a network of autonomous processes interacting through messages, evolving over time, and embedded in a shared computational environment. This chapter builds directly on that observation. Rather than \emph{discovering} agents as an emergent property of the infrastructure, we now \emph{design} them deliberately. The question is no longer whether agents exist on the Prolog Web, but how Prolog can be used to structure their behaviour, control their interaction, and integrate reasoning, learning, and tool use in a principled and inspectable way.


\section{Framing: why Prolog agents now}
\label{sec:framing-why-prolog-agents-now}

The classical Web is primarily a publishing and integration substrate for humans: documents, hyperlinks, and request--response endpoints. A plausible next phase -- one that is already taking shape -- is a Web populated by long-lived software processes that communicate, maintain state, and take actions over time. The technical heart of the shift is the change in the unit of composition, from pages and endpoints to interacting agents.

The term \emph{agentic} can sound like it merely denotes a new user-interface trend: chatty front ends, tool use, and LLM wrappers. In this chapter we use it in a stricter, architectural sense. An \emph{Agentic Web} is not defined by any particular interaction modality, but by the prevalence of autonomous, communicating processes as first-class web participants. By \emph{Agentic Prolog Web} we mean the fragment of that environment in which agents are implemented as Prolog processes with knowledge bases and reasoning procedures, exposed as web-native services and capable of interacting through message passing.

Historically, the Web has standardised naming, transport, and representation far more than it has standardised \emph{process structure}. We have excellent machinery for addressing and exchanging resources, but comparatively little is fixed about the lifecycle of long-lived services, the discipline of message-oriented interaction, the handling of partial failure, or the way stateful components compose. Early discussions of web agents within the Semantic Web community made this gap explicit: representation and inference mechanisms were advancing, but the process infrastructure in which agents would live and cooperate remained largely an application-level concern \citep{hendler2001agents}.

In practice, this means that agent behaviour tends to accrete in framework glue rather than being captured by a shared, inspectable execution model. The Prolog Web treats this gap as a design target. It keeps the Web's strengths -- protocols, naming, interoperability -- but adds an explicit execution model: nodes host long-lived actors with mailboxes and supervised lifecycles, and predicates can be exposed as message-driven services. This is why the chapter can move from ``agents as an emergent property'' to ``agents as a deliberate design'': the substrate already has the right shape for autonomous processes, and the remaining problem becomes behavioural engineering -- how to specify, constrain, explain, and govern what agents do.

That the architectural gap exists does not, by itself, explain why Prolog should be the language to fill it. Two contemporary pressures make the case.

The first is that agents need enforceable semantics. Large language models have made fluent natural-language interaction cheap and widely available, and they are already becoming an interface layer for many tools. However, fluency is not the same as behavioural reliability. Whenever an agent is expected to respect non-negotiable requirements
%-- permissions, invariants, resource budgets, safety policies, contractual obligations, or domain constraints -- 
we need a component whose behaviour has clear meaning and can be validated. Prolog offers precisely this: executable symbolic representations with inspectable derivations, well-defined failure modes, and a clear separation between what is \emph{entailed} and what is merely \emph{suggested}. In other words, Prolog is a natural place to put the parts of an agent that must be auditable, reproducible, and enforceable.

The second pressure is that the modern Web removes the old deployment bottleneck for symbolic AI. Historically, symbolic systems often remained laboratory artifacts: brittle deployments, heavy bespoke infrastructure, and limited interoperability. If an agent had to be installed, configured, and maintained as a special-purpose application, the economic and organisational cost could dominate the technical benefit. Standard protocols, cloud deployment practices, and ubiquitous client platforms change this equation. A Prolog engine can sit behind stable interfaces as a verifier, planner, policy engine, dialogue manager, or explanation service, and it can be composed with other services without losing its semantics.

The resulting positioning is not ``Prolog everywhere,'' but Prolog where semantic clarity and constraint satisfaction matter -- what we might call a \emph{behavioural backplane} for agentic systems: the component in which constraints, state, control, and explanations are represented as programs and can therefore supervise more fallible components. This positioning requires candour about Prolog's limitations. A steep learning curve, a fragmented implementation landscape, and limited mainstream ecosystem momentum all constrain adoption. The bet is that these costs are acceptable when the alternative -- building constraint enforcement and auditability ad hoc on top of a language that has neither -- is worse.

\medskip

\noindent The remainder of this chapter develops the argument in four steps. We begin with a compact reminder of Prolog's symbolic competence on a single machine: knowledge representation, systematic search, planning, meta-interpretation, explanation, and grammar-based language processing. We then lift these idioms onto the Agentic Web by treating Prolog engines as networked actors, so that predicates become message-driven services and proofs can span multiple nodes. Next, we introduce a neuro-symbolic layer in which large language models provide linguistic robustness while Web Prolog supplies the authoritative control and constraint envelope. Finally, we turn to human-facing user-interface agents, where dialogue is best understood as a control problem with states, repairs, and timeouts.


%\chapter{Prolog AI agents on the Agentic Web}
%\label{sec:prolog-ai-agents-on-the-agentic-web}
%
%Figure~\ref{fig:ai-and-prolog} sketches a popular retrospective view of Artificial Intelligence as a sequence of booms and busts: periods of intense optimism and investment followed by ``AI winters''. On this view, we are currently in a new boom whose scale and public visibility surpass earlier waves.
%
%\begin{figure}[h]
%\centering
%\includegraphics[width=11cm]{ai_booms.jpg}
%\caption{Cycle of AI booms and busts. Chart inspired by Yutaka Matsuo.}
%\label{fig:ai-and-prolog}
%\end{figure}
%
%\noindent Prolog and Lisp were prominent during earlier waves, especially the expert-systems era, whereas the current boom is driven primarily by data-intensive methods such as deep learning and large language models. This chapter argues that the story does not end there. As software becomes increasingly \emph{agentic}, the need for controllable behaviour and auditable reasoning re-emerges as a first-class engineering concern.
%
%We use the term \emph{Agentic Web} for the Web viewed as an environment populated by long-lived, interacting processes, and \emph{Agentic Prolog Web} for the fragment of that environment in which agents are implemented as Prolog processes with knowledge bases and reasoning procedures, exposed as web-native services and capable of interacting through message passing.
%
%Earlier in the book, and in particular in Chapter~\ref{sec:the-prolog-web} we argued that the Prolog Web already constitutes a multi-agent system in a minimal but precise sense: a network of autonomous processes interacting through messages, evolving over time, and embedded in a shared computational environment. This chapter builds directly on that observation. Rather than \emph{discovering} agents as an emergent property of the infrastructure, we now \emph{design} them deliberately. The question is no longer whether agents exist on the Prolog Web, but how Prolog can be used to structure their behaviour, control their interaction, and integrate reasoning, learning, and tool use in a principled and inspectable way.
%
%\section{Framing: why Prolog agents now}
%\label{sec:framing-why-prolog-agents-now}
%
%The present AI boom matters not only because models have improved, but because it is changing the \emph{unit} of software we build. Much of what is now marketed as ``AI'' is better understood as a shift from programs that merely transform inputs to outputs to programs that persist, maintain state, call tools, coordinate with other services, and take actions over time. In other words, we are increasingly deploying \emph{agents} rather than isolated functions.
%
%This shift makes a familiar symbolic-AI concern operational again: governing behaviour under non-negotiable constraints. A long-lived agent is not judged only by whether it produces plausible text. It is judged by whether it respects permissions, invariants, budgets, safety policies, and contractual obligations while operating in an environment with partial failure, latency, and conflicting demands. In such settings, ``good outputs most of the time'' is not an acceptable specification. The system needs a component that can represent what is allowed, decide what to do next, and justify why it did it.
%
%In that environment, the practical question for Prolog is not whether logic programming remains \emph{theoretically} relevant, but whether its strengths can be made \emph{operational}: deployable, composable, and inspectable inside modern software stacks. This chapter argues that the answer can be ``yes'' when Prolog is treated not as a monolithic application language but as a \emph{behavioural backplane} for agentic systems -- the place where constraints, policies, search, and explanations live as executable artifacts, and where more fallible components are invoked under explicit control.
%
%Two pressures shape the design space. First, agents need an enforceable semantics somewhere in the stack: a layer that can validate, refuse, and explain. Second, the modern Web makes specialised symbolic services economically viable: a Prolog engine can be deployed as a networked component behind stable interfaces, scaled and supervised with ordinary infrastructure, and composed with other services without giving up its meaning.
%
%\subsection{Agents need semantics}
%
%Large language models make fluent natural-language interaction cheap and widely available, and they are already becoming an interface layer for tools and services. However, fluency is not the same as behavioural reliability. When an agent must satisfy hard requirements -- permissions, invariants, resource budgets, safety policies, or domain constraints -- we need a component whose behaviour has clear meaning and whose outputs can be validated. Prolog offers precisely this: executable symbolic representations with inspectable derivations, explicit failure modes, and a sharp distinction between what is \emph{entailed} and what is merely \emph{suggested}. In practice, this makes Prolog a natural home for the parts of an agent that must be auditable, reproducible, and enforceable.
%
%\subsection{The Web removes the old bottleneck for symbolic AI}
%
%Historically, many symbolic systems remained costly to deploy: bespoke runtimes, brittle integration, and limited interoperability could dominate any technical gain. The modern Web changes this equation. Standard protocols, cloud deployment practices, and ubiquitous client platforms make it realistic to expose specialised reasoning components as networked services. In that setting, a Prolog engine becomes more valuable, not less: it can sit behind stable interfaces as a verifier, planner, policy engine, dialogue manager, or explanation service, and it can be composed with other components without losing its semantics.
%
%From this perspective, Prolog's future is best framed as a niche of high leverage: not ``Prolog everywhere'', but Prolog where semantic clarity and constraint satisfaction matter. The remainder of this chapter develops that stance concretely by first recalling Prolog's symbolic competence on a single machine, then lifting the same idioms onto the Agentic Prolog Web as message-driven services and interacting agents, and finally showing how such agents can wrap neural components inside an explicit control and constraint envelope.
%
%
%\section{Framing: why Prolog agents now}
%\label{sec:framing-why-prolog-agents-now}
%
%The present AI boom is not merely a breakthrough in pattern recognition. It is also a shift in how software is built and operated. Increasingly, programs are no longer conceived as static artifacts that transform inputs to outputs, but as long-lived processes that perceive, decide, act, and coordinate. In industry writing, this shift is often described as \emph{agentic}, and it motivates the framing adopted in this chapter.
%
%In such an environment, the practical question for Prolog is not whether logic programming remains \emph{theoretically} relevant. The question is whether Prolog's strengths can be made \emph{operational}: deployable at Internet scale, composable with mainstream infrastructure, and usable as a trustworthy substrate for agents that must act under constraints. This chapter argues that the answer can be ``yes'' when Prolog is treated not as a monolithic application language, but as a \emph{behavioural backplane} for agentic systems -- the component in which constraints, state, control, and explanations are represented as programs and can therefore supervise more fallible components.
%
%Two pressures shape the design space: the need for enforceable semantics inside agents, and the fact that the modern Web makes specialised symbolic services deployable and composable at scale.
%
%
%\subsection{Agents need semantics}
%
%Large language models have made fluent natural-language interaction cheap and widely available, and they are already becoming an interface layer for many tools. However, fluency is not the same as behavioural reliability. Whenever an agent is expected to respect non-negotiable requirements -- permissions, invariants, resource budgets, safety policies, contractual obligations, or domain constraints -- we need a component whose behaviour has clear meaning and can be validated. Prolog offers precisely this: executable symbolic representations with inspectable derivations, well-defined failure modes, and a clear separation between what is \emph{entailed} and what is merely \emph{suggested}. In other words, Prolog is a natural place to put the parts of an agent that must be auditable, reproducible, and enforceable.
%
%\subsection{The Web removes the old bottleneck for symbolic AI}
%
%Historically, symbolic AI systems often remained laboratory artifacts: brittle deployments, heavy bespoke infrastructure, and limited interoperability. If an agent had to be installed, configured, and maintained as a special-purpose application, the economic and organisational cost could dominate the technical benefit. The modern Web changes this equation. Standard protocols, cloud deployment practices, and ubiquitous client platforms make it realistic to deploy specialised reasoning services as networked components. In that setting, a Prolog engine becomes more valuable, not less: it can sit behind stable interfaces as a verifier, planner, policy engine, dialogue manager, or explanation service, and it can be composed with other services without losing its semantics.
%
%From this perspective, Prolog's future is best framed as a niche of high leverage: not ``Prolog everywhere'', but Prolog where semantic clarity and constraint satisfaction matter. Table~\ref{tab:swot-prolog-ai} summarises this positioning in the form of a SWOT matrix.
%
%\begin{table}[h]
%\centering
%\setlength{\tabcolsep}{6pt}
%\renewcommand{\arraystretch}{1.25}
%\begin{tabular}{|p{0.47\linewidth}|p{0.47\linewidth}|}
%\hline
%\textbf{Strengths} & \textbf{Weaknesses} \\
%\hline
%\textbf{Declarative semantics with executable meaning.} Programs express what holds, not only how to compute it. & \textbf{Steep conceptual learning curve.} Backtracking and control constructs such as the cut remain unintuitive to many programmers. \\
%\hline
%\textbf{Native support for search and constraints.} Planning, validation, and configuration are first-class concerns. & \textbf{Fragmentation across implementations.} Divergent extensions and libraries hinder portability and standardization. \\
%\hline
%\textbf{Explainability by construction.} Derivations, failures, and constraints can be inspected and justified. & \textbf{Limited mainstream ecosystem momentum.} Tooling, tutorials, and hiring pipelines lag behind dominant platforms. \\
%\hline
%\textbf{Metalanguage and interpreter-building strengths.} Well suited for reasoning about rules, languages, and models. & \textbf{Perceived unpredictability in performance.} Cost models are less obvious than in strictly deterministic languages. \\
%\hline
%\end{tabular}
%
%\vspace{0.8em}
%
%\begin{tabular}{|p{0.47\linewidth}|p{0.47\linewidth}|}
%\hline
%\textbf{Opportunities} & \textbf{Threats} \\
%\hline
%\textbf{AI as generator, Prolog as verifier.} LLMs propose solutions; Prolog checks constraints, invariants, and policies. & \textbf{Vendor-controlled agent platforms.} Closed ecosystems may marginalize open symbolic reasoning engines. \\
%\hline
%\textbf{Agent-based systems with explicit reasoning cores.} Renewed interest in symbolic planning and rule-based decision making. & \textbf{Natural-language programming obscuring semantics.} ``Good enough'' behaviour may be preferred over formally grounded logic. \\
%\hline
%\textbf{Executable policies and compliance rules.} Growing demand for auditable, machine-checkable specifications. & \textbf{Integration gravity toward mainstream runtimes.} WASM, JVM, and cloud-native stacks may dominate interfaces and tooling. \\
%\hline
%\textbf{Hybrid neuro-symbolic architectures.} Statistical models combined with logical constraints and explanations. & \textbf{Security concerns around dynamic code.} Systems lacking clear sandboxing and resource control risk exclusion. \\
%\hline
%\end{tabular}
%\caption{SWOT matrix: Prolog in an AI-driven software landscape.}
%\label{tab:swot-prolog-ai}
%\end{table}
%
%
%\subsection{From the Web of documents to the Agentic Web}
%\label{sec:from-documents-to-agentic-web}
%
%The classical Web is primarily a publishing and integration substrate for humans: documents, hyperlinks, and request--response endpoints. The Agentic Web is a plausible next phase: the same network substrate, but populated by long-lived software processes that communicate, maintain state, and take actions over time. The technical heart of the shift is the change in the unit of composition, from pages and endpoints to interacting agents.
%
%\subsubsection{Agents as a web-level unit of composition}
%The term \emph{agentic} can sound like it merely denotes a new user-interface trend: chatty front ends, tool use, and LLM wrappers. In this chapter we use it in a stricter, architectural sense. An Agentic Web is not defined by any particular interaction modality, but by the prevalence of autonomous, communicating processes as first-class web participants.
%
%\subsubsection{What the classical Web does not standardise}
%Historically, the Web has standardised naming, transport, and representation far more than it has standardised \emph{process structure}. We have excellent machinery for addressing and exchanging resources, but comparatively little is fixed about the lifecycle of long-lived services, the discipline of message-oriented interaction, the handling of partial failure, or the way stateful components compose. Early discussions of web agents within the Semantic Web community made this gap explicit: representation and inference mechanisms were advancing, but the process infrastructure in which agents would live and cooperate remained largely an application-level concern \citep{hendler2001agents}.
%
%In practice, this means that agent behaviour tends to accrete in framework glue rather than being captured by a shared, inspectable execution model. The Prolog Web, and in particular the Agentic Prolog Web, treats this gap as a design target. It keeps the Web's strengths -- protocols, naming, interoperability -- but adds an explicit execution model: nodes host long-lived actors with mailboxes and supervised lifecycles, and predicates can be exposed as message-driven services. This is why the chapter can move from ``agents as an emergent property'' to ``agents as a deliberate design'': the substrate already has the right shape for autonomous processes, and the remaining problem becomes behavioural engineering -- how to specify, constrain, explain, and govern what agents do.
%
%\medskip
%
%\noindent The remainder of this chapter develops the argument in four steps. We begin with a compact reminder of Prolog's symbolic competence on a single machine: knowledge representation, systematic search, planning, meta-interpretation, explanation, and grammar-based language processing. We then lift these idioms onto the Agentic Web by treating Prolog engines as networked actors, so that predicates become message-driven services and proofs can span multiple nodes. Next, we introduce a neuro-symbolic layer in which large language models provide linguistic robustness while Web Prolog supplies the authoritative control and constraint envelope. Finally, we turn to human-facing user-interface agents, where dialogue is best understood as a control problem with states, repairs, and timeouts.

\section{Symbolic competence in one engine}
\label{sec:symbolic-prolog}

Before we distribute Prolog across the Agentic Web, it is useful to recall what a single Prolog engine already provides as an AI substrate. The point of this section is not to teach Prolog, but to fix a compact capability map: a handful of idioms that repeatedly recur in symbolic AI, and that will later reappear as message-driven services and interacting agents.

\subsection{Knowledge representation and inference}

A Prolog program can function simultaneously as a declarative knowledge base and as an executable inference procedure. A canonical example is a small relational theory of kinship, where recursion expresses transitive closure:

\begin{verbatim}
ancestor_descendant(X, Y) :- parent_child(X, Y).
ancestor_descendant(X, Z) :-
    parent_child(X, Y), ancestor_descendant(Y, Z).

parent_child(X, Y) :- mother_child(X, Y).
parent_child(X, Y) :- father_child(X, Y).

mother_child(trude, sally).

father_child(tom, sally).
father_child(tom, erica).
father_child(joe, tom).
\end{verbatim}

\noindent A query such as \texttt{?-ancestor\_descendant(joe,\,Who)} asks for logical consequences of the clauses and enumerates answers via backtracking. This combination -- explicit relational structure plus a proof procedure -- is the basic engine behind many symbolic systems, and in Prolog the proof procedure is built in.

\subsection{Combinatorial search: a classic N-Queens program}
\label{sec:symbolic-nqueens}

Many agent tasks reduce to combinatorial exploration: scheduling, placement, configuration, and other problems where one must search a large discrete space under hard constraints. Even in plain ISO-style Prolog, such tasks can often be expressed compactly by representing partial solutions as terms and letting the engine explore alternatives under backtracking, with carefully chosen data representations to prune the search early.

A traditional Prolog solution to the N-Queens puzzle illustrates the style. The program below, due to Richard \citet{okeefe1990craft}, encodes the board as a structured set of constraints and constructs a consistent placement row by row:

\begin{verbatim}
queens(N, Queens) :-
    length(Queens, N),
    board(Queens, Board, 0, N, _, _),
    queens(Board, 0, Queens).

board([], [], N, N, _, _).
board([_|Queens], [Col-Vars|Board], Col0, N, [_|VR], VC) :-
    Col is Col0+1,
    functor(Vars, f, N),
    constraints(N, Vars, VR, VC),
    board(Queens, Board, Col, N, VR, [_|VC]).

constraints(0, _, _, _) :- !.
constraints(N, Row, [R|Rs], [C|Cs]) :-
    arg(N, Row, R-C),
    M is N-1,
    constraints(M, Row, Rs, Cs).

queens([], _, []).
queens([C|Cs], Row0, [Col|Solution]) :-
    Row is Row0+1,
    select(Col-Vars, [C|Cs], Board),
    arg(Row, Vars, Row-Row),
    queens(Board, Row, Solution).
\end{verbatim}

\noindent The point of including this program in a Web Prolog context is not to claim that a minimal, ISO-oriented profile is the best vehicle for large-scale constraint programming. Rather, the example marks a capability boundary: even without specialised constraint libraries, Prolog can express structured search problems declaratively and solve them by systematic exploration. Where a Prolog system does provide additional facilities such as finite-domain constraints, tabling, or specialised propagators, the same modelling stance can be retained while the engine becomes substantially more efficient. In the Web Prolog profile used throughout this book, however, we restrict ourselves to ISO-level assumptions unless otherwise stated.

\subsection{Planning as search in a state-transition model}

Many agent tasks are naturally expressed as planning: find a sequence of actions that transforms an initial state into a goal state. A small meta-program suffices when actions are represented as declarative transition rules:

\begin{verbatim}
plan(P, P, []) .
plan(P0, Pf, [A|As]) :-
    step(P0, P1, A),
    plan(P1, Pf, As).

step(start, position(monkey, bananas, chair), start).
step(position(Y,Y,Y), final, climb).
step(position(X,Y,X), position(Y,Y,Y), carry(X,Y)).
step(position(X,Y,Z), position(Z,Y,Z), walk(X,Z)).
\end{verbatim}

\noindent Here the entire world state is a first-order term, and each \texttt{step/3} clause specifies an action's preconditions and effects. Prolog's search then discovers action sequences by exploring transition alternatives under backtracking.

\subsection{Meta-interpreters as a control and explanation lever}
\label{sec:basic-meta-interpreter}

A distinctive Prolog idiom is to treat programs as data and write interpreters for restricted logics, alternative control regimes, or explanation-producing variants. A minimal meta-interpreter mirrors Prolog's own proof search:

\begin{verbatim}
prove(true) :- !.
prove((A,B)) :- !, prove(A), prove(B).
prove(A) :- clause(A,B), prove(B).
\end{verbatim}

\noindent Small changes to such an interpreter can add substantial functionality without changing the object-level knowledge base. For example, adding a second argument allows the interpreter to construct a proof object -- a structured explanation that records which rules were used:

\begin{verbatim}
prove(true, true) :- !.
prove((A,B), (PA,PB)) :- !, prove(A, PA), prove(B, PB).
prove(A, A/Proof) :- clause(A,B), prove(B, Proof).
\end{verbatim}

\noindent Proof terms of this kind are evidence-carrying data structures that can be inspected, summarised, or shown to a user. In Section~\ref{sec:prolog-web-agents} we will exploit this idiom across nodes, so that a single explanation can span distributed computation.

\subsection{An expert system with a query-the-user facility}
\label{sec:query-the-user}

Expert systems were a flagship application of symbolic AI during the 1980s. They did not become a lasting mainstream paradigm, but the core pattern remains instructive, especially in Prolog: it shows how deduction can be interleaved with interaction in a disciplined way. 

An expert system typically consists of a knowledge base plus an inference engine. A practical system must also be able to acquire facts that are not already stored. In Prolog, a simple way to realise this is to extend the meta-interpreter from Section~\ref{sec:basic-meta-interpreter} with a single additional clause that queries the user whenever a goal is declared \texttt{askable}:

%\begin{verbatim}
%prove(A) :-
%    askable(A, Question),
%    writeln(Question),
%    read(Answer),
%    Answer == yes.
%\end{verbatim}

\begin{verbatim}
prove(A) :- askable(A, Q), writeln(Q), read(T), T == yes.
\end{verbatim}

\noindent The accompanying knowledge base supplies both domain rules and a set of \texttt{askable/2} declarations that map predicates to human-readable prompts:

\begin{verbatim}
good_pet(X) :- bird(X), small(X).
good_pet(X) :- cuddly(X), yellow(X).

bird(X) :- has_feathers(X), tweets(X).

askable(tweets(_), 'Does it tweet?').
askable(small(_), 'Is it small?').
askable(cuddly(_), 'Is it cuddly?').
askable(has_feathers(_), 'Does it have feathers?').
askable(yellow(_), 'Is it yellow?').
\end{verbatim}

\noindent When \texttt{prove/1} encounters an askable subgoal, it suspends deduction, poses the associated question, reads the reply, and then resumes the derivation. For example:

\begin{verbatim}
?- prove(good_pet(tweety)).
Does it have feathers?
|: yes.
Does it tweet?
|: yes.
Is it small?
|: no.
Is it cuddly?
|: yes.
Is it yellow?
|: yes.
true.
?-
\end{verbatim}

\noindent The important point is structural: proof search determines which information is needed when. The resulting interaction is mixed-initiative in a precise sense. The engine advances the derivation whenever its knowledge suffices, and it consults the user exactly at the proof obligations that require additional axioms. The \texttt{askable} declarations therefore act both as an interaction interface and as a lightweight user model: they mark which facts are expected to come from dialogue rather than from stored knowledge.

This is, of course, a toy shell. In more realistic systems one would add typed answers, richer question forms, explanation facilities, and a persistent dialogue state. For a fuller treatment of Prolog-based expert systems, see ~\cite{Merritt1989} and \cite{sterling1994art}.

\subsection{Parsing as deduction: DCGs}

Definite Clause Grammars (DCGs) turn parsing into a form of proof search: a sentence is grammatical exactly when it can be derived from the grammar rules. DCGs are also reversible, so the same grammar can be used for generation. Operationally, naive top-down execution makes left recursion problematic, but in tabled Prolog systems this restriction largely disappears: tabling turns parsing into dynamic programming while preserving a declarative rule notation. (Tabling is a system feature rather than part of the ISO core, so it should be read here as an optional acceleration of the same modelling stance.)

\subsection{Summary and handover}

The examples above form a compact capability map. A single Prolog engine already supports (i) explicit relational knowledge representation, (ii) systematic search with explicit structure and early pruning, (iii) planning in state-transition models, (iv) meta-level control and explanation via interpreters that construct proof objects, (v) interactive acquisition of missing facts, and (vi) grammar-based language processing as deduction.

The next section lifts these idioms onto the Agentic Web. Instead of treating predicates as internal procedures, we treat them as services: message-driven actors with interfaces, state, and failure modes. The result is an Agentic Prolog Web in which symbolic reasoning is not only executable, but deployable, composable, and inspectable across the network.

\section{Prolog AI agents on the Agentic Prolog Web}
\label{sec:prolog-web-agents}

The central claim of this section is that the idioms (i)--(vi) do not need to remain local. When Prolog engines are exposed as web-native actors, the same programs can be deployed as distributed services and composed into interacting agents.

The Agentic Prolog Web is therefore not a new symbolic formalism. It is an execution and deployment stance: treat Prolog engines as networked processes with mailboxes, state, and interfaces, and let predicates become message-driven services. In this setting, what used to be a local \texttt{read/1} becomes a \texttt{prompt/2} message expecting a response; what used to be a local meta-interpreter becomes a mobile proof-producing component; and what used to be an internal predicate call may become a remote call whose provenance is part of the explanation.

We develop the idea through four ingredients: an expert-system agent that acquires missing facts as a distributed dialogue; a proof-producing meta-interpreter whose explanations span multiple nodes; a positioning of Web Prolog relative to earlier Prolog-based agent frameworks; and a brief discussion of scalability and interoperability on the Web.

\subsection{An expert-system agent and a simulation}
\label{sec:expert-system-agent}

Classic expert systems faced two persistent bottlenecks: the difficulty of extracting knowledge from humans and the lack of a ubiquitous infrastructure for deployment. A web-native actor substrate addresses both. By treating the Web as the infrastructure, symbolic agents can be deployed as services; and by treating interaction as message passing, knowledge acquisition becomes an incremental, mixed-initiative dialogue rather than a local read/write loop.

To adapt the askable meta-interpreter pattern from Section~\ref{sec:query-the-user} to an agent setting, we replace terminal I/O with an interaction predicate. Concretely, the clause that previously wrote a question and read an answer becomes:

\begin{verbatim}
prove(A) :- askable(A, Q), input(Q, T), T == yes.
\end{verbatim}

\noindent The purpose of \texttt{input/2} is to decouple deduction from any particular user interface. In a local setting, \texttt{input/2} might still be implemented on top of \texttt{read/1}. As we saw in Chapter~\ref{sec:the-prolog-web}, in the Prolog Web it is implemented by sending a \texttt{prompt} to some conversational partner and waiting for a reply via \texttt{respond/2}.

We can now package the expert system as an expert-system agent by spawning a remote toplevel actor and preparing a call to \texttt{prove/1}:

\begin{verbatim}
expert_ask(Query, Pid, Options) :-
    toplevel_spawn(Pid, Options),
    toplevel_call(Pid, prove(Query), [
        limit(1)
      | Options
    ]).
\end{verbatim}

\noindent The \texttt{limit(1)} option is not an arbitrary restriction. It enforces a simple turn-taking discipline where each remote call either produces a single \texttt{prompt} or a single terminal outcome (\texttt{success} or \texttt{failure}).\footnote{Without such a bound, solution enumeration could interleave multiple answers with interaction steps, complicating the conversational protocol and making it unclear which prompt belongs to which branch of search.}

To demonstrate the pattern, we construct a small simulation in which a second agent -- the shell process in fact -- plays the role of the user. The simulation answers questions stochastically, illustrating that the expert system does not control the dialogue by a script; rather, the proof search itself determines which information must be acquired next.

% loading it with the required predicates and facts, and issuing a query:

\begin{verbatim}
simulation(Options) :-
    random_member(Name, [tweety,pingu]),
    format('A1: Is ~w a good pet?~n', [Name]),
    expert_ask(good_pet(Name), Pid, Options),
    interact(Pid, Name).

interact(Pid, Name) :-
    receive({
        failure(Pid) ->
            format('A2: No, ~w is not a good pet.~n', [Name]);
        success(Pid, _, _) ->
            format('A2: Yes, ~w is a good pet.~n', [Name]);
        prompt(Pid, Question) ->
            format('A2: ~w~n', [Question]),
            random_member(Answer, [yes,no]),
            respond(Pid, Answer),
            format('A1: ~w~n', [Answer]),
            interact(Pid, Name)
    }).
\end{verbatim}

\noindent One may run the simulation, for example, by spawning the expert-system agent on a remote node, loading the relevant predicates, and providing an initial fact base:

\begin{verbatim}
?- simulation([
       node('http://n9.org'),
       load_predicates([prove/1]),
       load_list([yellow(tweety)])
   ]).
A1: Is tweety a good pet?
A2: Does it have feathers?
A1: yes
A2: Does it tweet?
A1: yes
A2: Is it small?
A1: no
A2: Is it cuddly?
A1: yes
A2: Yes, tweety is a good pet.
true.
?-
\end{verbatim}

\noindent The architectural point is that the reasoning process is encapsulated as a mobile, interactive component: an actor whose control flow is governed by logical necessity rather than by a fixed dialogue script. The Web substrate supplies distribution and composability, while Prolog supplies the proof structure that drives which questions must be asked.

\subsection{Explainable AI across nodes: distributed proof trees}
\label{sec:distributed-proof-trees}

A defining advantage of logic-based, symbolic AI is transparency: every conclusion rests on rules and facts, and an explanation can be constructed as a proof object. On the Agentic Prolog Web, explanations must survive distribution. If a subgoal is proved on another node, its justification should be folded into the same overall proof structure rather than relegated to an opaque remote call.

A compact meta-interpreter can achieve this by treating remote calls as proof-producing sub-derivations. When the interpreter encounters a remote goal, it requests a proof-producing interpreter on the remote node and continues constructing the proof there:

\begin{verbatim}
prove(true, true) :- !.
prove(rpc(URI,A), Proof) :- !,
    prove(rpc(URI,A,[]), Proof).
prove(rpc(URI,A,Options), Query@URI/Proof) :- !,
    rpc(URI, prove(A, Query/Proof), [
        load_predicates([prove/2])
      | Options
    ]).
prove((A, B), (ProofA, ProofB)) :- !,
    prove(A, ProofA),
    prove(B, ProofB).
prove(A, A/Proof) :-
    clause(A, B),
    prove(B, Proof).
\end{verbatim}

\noindent The key idea is that remote computation is not treated as a black box. Instead, the same explanation discipline is applied at each node, and remote contributions are marked in the resulting proof term.

Suppose the local node contains the rule \texttt{mortal(X)\,:-\,human(X)} and also defines \texttt{human(X)} by consulting a remote node. If the fact \texttt{human(plato)} resides on \texttt{n1.org}, we obtain a proof that records the distribution boundary:

\begin{verbatim}
?- prove(mortal(plato), Proof).
Proof = mortal(plato)/
          (human(plato)/
             (human(plato)@'http://n1.org'/
                true)).
?-
\end{verbatim}

\noindent Proof-producing interpreters can of course generate very large derivations. On the Prolog Web this is best treated as a controllable engineering choice rather than an unavoidable pathology: one may bound proof depth, abstract away local details and retain only a node-level skeleton, cache recurring subproofs, or compress proof objects into summaries suitable for user-facing explanations. The benefit is that distribution does not force a loss of transparency. Explanations remain first-class data, even when computation spans multiple machines.

\subsection{Parsing as remote deduction}
\label{sec:remote-parsing}

A client may offload parsing to a node that hosts a grammar, without embedding that grammar locally:

\begin{verbatim}
?- rpc('http://n6.org', phrase(s(PT), [joe,saw,a,telescope])).
PT = s(np(pn(joe)), vp(v(saw), np(det(a), n(telescope)))).
?-
\end{verbatim}

\noindent This vignette illustrates the broader pattern: symbolic competencies that are expensive to ship as libraries or models can be deployed as web-native agents and accessed by clients through uniform remote calls.

\subsection{Implementing agent frameworks in Web Prolog}
\label{sec:agent-frameworks-on-web-prolog}

Logic-based agent languages and frameworks have been implemented in Prolog for decades. Systems such as DALI \citep{costantini2004dali} and LPS \citep{kowalski2015reactive} show that one can retain a Horn-clause core while adding structures for events, actions, reactive rules, and persistent state. Their appeal lies in combining logic programming virtues -- declarative clarity, inspectability, and verifiability -- with practical mechanisms for autonomy and interaction.

Web Prolog does not present itself as yet another high-level cognitive agent language with built-in planners, ontologies, or a fixed BDI-style architecture. Instead, it provides a lower-level actor substrate in which the primary abstraction is a process with a private database and a mailbox, encapsulated behind a process identifier and supervised by reliability patterns. In this respect it is closer to an Erlang-inspired concurrent logic substrate than to a monolithic agent language.

Precisely because the substrate is low-level and concrete, many classic agent-language ideas can be reconstructed. A receive loop that pattern-matches on incoming messages, together with an explicit representation of pending goals, already suffices to encode reactive rules of the form ``on event then action'' or ``on event and if condition, then action'' as ordinary Prolog clauses over the mailbox and the private database. In statechart actors, this becomes even more structured, allowing rich reactive and mixed-initiative behaviours to be expressed as statechart specifications interpreted by Web Prolog rather than as ad hoc control code.

Seen from this perspective, Web Prolog is again best understood as an enabler rather than a competitor. It supplies the web-native execution fabric on top of which more specialised agent frameworks can be layered. LPS-style engines, DALI-like event-driven agents, and more elaborate frameworks can all be realised as conservative extensions that treat Web Prolog actors as their concrete embodiment on the open Web.

\subsection{Scalability and interoperability on the Web}
\label{sec:scalability-and-interoperability}

Two traits of the Web are decisive for symbolic agents: elastic scalability and protocol-level interoperability.

\subsubsection{Scalability}

Web Prolog agents are processes that can be spawned, suspended, or terminated on demand. Because interaction can be mediated by standard web protocols, off-the-shelf deployment practices apply: container orchestration, replication, load balancing, and ordinary observability tooling can be used to scale symbolic services without inventing bespoke infrastructure. Importantly, the actor model supplies natural failure boundaries: one can supervise, restart, and isolate components at the granularity of individual agents.

\subsubsection{Interoperability}

The Web rewards components that speak stable protocols. On the Agentic Prolog Web, terms travel as structured data, which makes it natural to expose symbolic services as typed message interfaces and to interoperate with non-Prolog components through canonical serialisations. Conversely, Prolog agents can call REST services, interact with browser clients, or orchestrate external systems, while keeping their internal reasoning symbolic and inspectable.

\medskip

\noindent Distributed symbolic agents already provide transparency and coordination at Internet scale, but many deployed agents are not purely agent-to-agent components. They are hybrid systems that must incorporate neural components for language, perception, and retrieval, while keeping authority in a checkable substrate. The next section therefore introduces neuro-symbolic agents: architectures in which Web Prolog acts as control and constraint envelope, and LLMs act as pluggable language modules.

\section{Neuro-symbolic AI agents}
\label{sec:neuro-symbolic}

\begin{quotation}
\noindent To build a robust, knowledge-driven approach to AI we must have the machinery of symbol manipulation in our toolkit. Too much useful knowledge is abstract to proceed without tools that represent and manipulate abstraction, and to date, the only known machinery that can manipulate such abstract knowledge reliably is the apparatus of symbol manipulation.
\flushright \textit{Marcus and Davis, 2019}
\vspace{0.5cm}
\end{quotation}

\noindent Large language models are powerful generators of fluent text, and that fluency is precisely why they are not reliability-first systems. When facts, constraints, obligations, or provenance are missing from the prompt, a model can still produce a plausible continuation, and the result may read as confident even when it is wrong. In low-stakes settings this is often acceptable and sometimes useful, but in settings where actions are irreversible or correctness is externally defined (healthcare, compliance, infrastructure automation), we need architectures that can enforce invariants, expose assumptions, and fail safely. Neuro-symbolic agents address this by combining neural components for perception and language with symbolic components for rules, search, and verification, so that generative flexibility is grounded in checkable structure rather than treated as the final authority.

Sections~\ref{sec:symbolic-prolog} and \ref{sec:prolog-web-agents} established, respectively, the symbolic competence of pure Prolog programs and their web-native, actor-oriented execution model. Building on that foundation, this section shows how neural components (LLMs) and symbolic components (Web Prolog) can be composed into coherent agents, and why Web Prolog is a natural host language for such hybrids when the agent is expected to be concurrent, auditable, and distributed.

The landscape of agentic and neuro-symbolic systems is evolving so quickly that any classification of ``current best practice'' should be read as provisional. Accordingly, this section focuses on a small set of relatively stable architectural ideas -- control (statecharts), constraints (symbolic rules), and supervision (actors) -- that remain valuable even as specific tools and model capabilities change.

A further reason to be modest about predictions is that the boundary between description and implementation is itself shifting. It is not implausible that, at some point, one will be able to feed the text of this book -- specifications, examples, protocol sketches, and semantic constraints -- to a capable LLM and obtain a working implementation of a Prolog node in return. If that happens, the purpose of the architectural material in this chapter is not undermined; rather, it becomes even more important to state the intended semantics clearly, to separate invariants from conventions, and to make the control and safety envelope explicit. Faster code generation increases the value of precise behavioural specifications: it lets us generate implementations more quickly while still knowing what, exactly, they are supposed to implement.

\subsection{Why hybridize?}
\label{sec:why-hybridize}

Neural models excel at perception, linguistic robustness, and learning from heterogeneous data, but they do not natively preserve invariants that a symbolic program can enforce: deductive consistency, traceability, structured explainability, and predictable failure modes. Hybridization therefore aims to combine (i) neural flexibility for perception and language understanding, (ii) symbolic rigour for logic, planning, and verification, and (iii) actor concurrency for scalable orchestration across the Web. The goal is not to replace neural methods, but to embed them in an architecture where the non-negotiables of a domain can be represented and checked continuously.

Table~\ref{tab:prolog-vs-llm-posture} contrasts the two paradigms at the level of engineering posture: what kind of semantics they offer, how they fail, and what can (or cannot) be guaranteed. It is meant to clarify why hybrid designs often assign authority to the symbolic layer even when language is handled by an LLM.

\begin{table}[ht]
\centering
\setlength{\tabcolsep}{6pt}
\renewcommand{\arraystretch}{1.25}
\begin{tabular}{|p{0.17\linewidth}|p{0.36\linewidth}|p{0.36\linewidth}|}
\hline
\textbf{Posture} & \textbf{Prolog-based reasoning systems} & \textbf{LLM-only systems} \\
\hline
\textbf{Semantics and guarantees} &
Precise operational meaning; behaviour can be constrained by rules, proof procedures, and resource bounds. &
Implicit and statistical; hard to guarantee correctness beyond narrow, externally enforced constraints. \\
\hline
\textbf{Correctness posture} &
Designed to be checked: constraints, invariants, and entailment can be validated and reproduced. &
Designed to be plausible: can be accurate, but may also hallucinate or drift without reliable internal validity checks. \\
\hline
\textbf{Explainability} &
Can return structured reasons (derivations, rule firings, failed constraints) and support audit trails. &
Explanations are generated narratives; may be persuasive but are not inherently evidence-carrying. \\
\hline
\textbf{Failure modes} &
Typically explicit: non-termination, incompleteness from control choices, or constraint inconsistency; failures can be instrumented. &
Often implicit: silent errors, confident wrong answers, instruction-following failures, and sensitivity to prompt phrasing. \\
\hline
\textbf{Data and learning} &
Knowledge is explicit and editable; learning is external (rule authoring, induction, or updates to facts/rules). &
Knowledge is latent in parameters; updating behaviour often requires retraining, fine-tuning, or additional scaffolding. \\
\hline
\textbf{Integration shape} &
Natural fit as a verifier, planner, policy engine, or reasoning service; can sit behind APIs with stable semantics. &
Natural fit as a generator and interface layer; needs surrounding machinery for validation, safety, and state consistency. \\
\hline
\textbf{Strength under constraints} &
Strong when requirements must be satisfied (safety, compliance, configuration, scheduling). &
Strong when requirements are soft and language-driven (drafting, summarizing, proposing options). \\
\hline
\textbf{Cost model} &
Predictable per inference step; performance depends on search space and control but can often be bounded. &
Token- and model-size-driven; costs can be high and scale with context length; behaviour may vary across model versions. \\
\hline
\end{tabular}
\caption{Contrasting Prolog-based reasoning systems with LLM-only systems: engineering posture and guarantees.}
\label{tab:prolog-vs-llm-posture}
\end{table}

%Table~\ref{tab:prolog-vs-llm-strengths} is a more pragmatic complement: it summarizes the kinds of tasks and evaluation criteria under which each approach tends to excel. Its role is not to declare a winner, but to make explicit where hybrid architectures typically place \emph{decision authority} versus \emph{language capability}.
%
%\begin{table}[ht]
%\centering
%\setlength{\tabcolsep}{6pt}
%\renewcommand{\arraystretch}{1.25}
%\begin{tabular}{|p{0.47\linewidth}|p{0.47\linewidth}|}
%\hline
%\textbf{Where Prolog tends to win} & \textbf{Where LLM-only tends to win} \\
%\hline
%\textbf{Verification and enforcement.} Checking constraints, invariants, permissions, and policy; producing counterexamples and minimal explanations. &
%\textbf{Fluent generation.} Drafting text, code sketches, test ideas, documentation, and interactive dialogue. \\
%\hline
%\textbf{Structured search.} Planning and combinatorial exploration with pruning, tabling, and constraint propagation. &
%\textbf{Heuristic exploration.} Rapidly proposing candidate solutions and creative alternatives without explicit search machinery. \\
%\hline
%\textbf{Auditability.} Evidence-carrying answers and reproducible reasoning steps. &
%\textbf{Adaptability.} Handling messy inputs, vague intent, and heterogeneous knowledge expressed in natural language. \\
%\hline
%\textbf{Stable semantics over time.} Given the same program and facts, results are stable (modulo control choices). &
%\textbf{Rapid improvement curve.} Capability can jump with new model releases, often without changes to user-facing prompts. \\
%\hline
%\end{tabular}
%\caption{Contrasting Prolog-based reasoning systems with LLM-only systems: typical strengths in practice.}
%\label{tab:prolog-vs-llm-strengths}
%\end{table}

\noindent A useful way to read the contrast between Prolog and LLM-only systems is in terms of where each one tends to carry engineering weight in deployed agents. Prolog tends to dominate where \emph{correctness is externally defined} and must be enforced rather than narrated. This includes verification and enforcement tasks such as checking constraints, invariants, permissions, and policy, where the system should be able to produce counterexamples or minimal justifications when a requirement is violated. It also includes structured search problems such as planning and combinatorial exploration, where explicit pruning, tabling, or constraint propagation can turn a large discrete space into an executable plan. Finally, Prolog retains a distinctive advantage whenever auditability is a first-class requirement: it can return evidence-carrying artefacts (derivations, failed constraints, proof objects) and support reproducible reasoning steps, and its meaning is comparatively stable over time in the sense that the same program and facts yield the same logical outcomes (modulo control choices).

LLM-only systems, by contrast, tend to dominate where the objective is fundamentally linguistic or heuristic rather than deductive. They excel at fluent generation: drafting text, code sketches, test ideas, documentation, and maintaining interactive dialogue that can tolerate underspecification. They also excel at heuristic exploration: rapidly proposing candidate solutions and creative alternatives without explicit search machinery. Their robustness to messy input is often their most practical advantage: they can accommodate vague intent and heterogeneous knowledge expressed in natural language where a symbolic interface would require careful schema design. Finally, their improvement curve can be unusually steep: capability can jump with new model releases even when the surrounding system and prompts remain unchanged. For hybrid architectures, the design takeaway is not that one side should replace the other, but that \emph{decision authority} is typically placed with the symbolic layer when constraints, audit, and stable meaning matter, while \emph{language capability} and heuristic proposal-making are delegated to the neural layer when the problem is communicative, underspecified, or primarily about drafting and interpretation.


\subsection{What is neuro-symbolic AI?}
\label{sec:what-is-neuro-symbolic}

A neuro-symbolic system is any computational architecture in which neural and symbolic modules exchange intermediate representations at run time under the control of an agent framework. The coupling may be loose (black-box tool calls with structured I/O), semi-tight (shared constraints, schemas, or retrieval paired with symbolic checks), or tight (a single differentiable pipeline in which parts of the symbolic procedure are compiled into a learnable computation graph). These three modes correspond to three engineering postures: orchestration, co-representation, and end-to-end training.

\subsection{Architectural patterns for hybrid agents}
\label{sec:patterns-hybrid-agents}

Four patterns recur in contemporary neuro-symbolic agent design.

\subsubsection{Reflection}

An LLM invokes external tools for objective sub-tasks such as retrieval, arithmetic, symbolic query answering, or constraint checking. Instead of fabricating a result, the neural component delegates, and the returned value constrains the next generation step. In practice, this is the simplest way to reduce factual drift and to obtain verifiable intermediate results.

\subsubsection{Planning with verification}

The LLM proposes an explicit action sequence while the symbolic layer checks preconditions, postconditions, and resource constraints before execution. The LLM is treated as a planner and summarizer, while Web Prolog plays the role of verifier and dispatcher. This pattern is especially useful when actions are partially ordered, when actions can fail, or when the environment imposes hard constraints that must never be violated.

\subsubsection{Multi-agent decomposition}

Sub-goals are assigned to cooperating actors so that each specialist works locally while the overall system remains concurrent and robust to partial failure. The neural component may participate as one actor among others, or it may be called by several actors as a shared service. In both cases, actor boundaries provide a natural granularity for supervision, timeouts, and audit trails.

\subsubsection{Finite-state or statechart cage}

The LLM is embedded inside a deterministic control wrapper that constrains what can happen next. The wrapper can accept, reject, or repair model outputs, and it can route the dialogue through explicit states such as \texttt{collect}, \texttt{confirm}, and \texttt{abort}. In Web Prolog this wrapper can be implemented directly as a statechart actor, giving the system a concrete behavioural skeleton rather than an implicit, prompt-only one.

A closely related engineering stance is articulated by Adam Charlson in his 2025 talk on building multi-agent systems with finite state machines. There, finite state machines (more precisely, statecharts) are presented as a governance layer around LLM-driven systems. The core idea matches the cage pattern used in this chapter: the LLM may propose actions or emit events, but an external statechart remains authoritative, enforcing valid transitions via guards and refusing to act on unmodelled events (which are instead logged for remediation). In this sense, state machines supply predictability, observability, and control that unconstrained text-generation loops typically lack, while the LLM contributes flexible language and task-specific heuristics inside the supervised control envelope.\footnote{This talk is available at \url{https://youtu.be/OD13PiXw60o}.}

\subsection{Web Prolog as symbolic backend and agent layer}
\label{sec:web-prolog-as-backend}

Python dominates ML practice for good reasons, and the point of this chapter is not to argue against Python. The point is that a hybrid agent needs more than tensor libraries: it needs an orchestration substrate with first-class support for search, symbolic constraints, concurrency, supervision, and network-transparent messaging. Web Prolog places these capabilities in the language model itself.

Web Prolog offers built-in backtracking and nondeterminism, which lets an agent represent alternatives and explore them under logical control rather than by ad hoc sampling. It supports guarded message reception and mailbox-driven coordination, which makes it natural to express reactive, concurrent control flows without collapsing everything into a single event loop. Its actor primitives and supervision patterns provide a principled way to handle partial failure, retries, and escalation. Finally, because messaging over the Web is part of the language-level model, the boundary between local and remote computation can be treated uniformly, which matters when a hybrid agent is distributed across nodes and services.

The result is a single substrate in which symbolic reasoning, orchestration, supervision, and network transparency are expressed in one semantics. Neural components then become pluggable co-processors rather than the place where control logic silently accumulates.

\subsubsection{Normative reasoning as executable policy}

Legal, ethical, and organisational constraints can be modelled as rule sets in Prolog, including defeasible rules, exceptions, and priority schemes. This makes it possible to treat compliance not as post-hoc documentation but as executable policy: before an action is taken, the agent can attempt to prove that it is permitted under the current facts, and if no proof exists it can either refuse, ask for missing information, or escalate to a human.

Even toy examples such as Asimov's laws are useful here, not because they are realistic, but because they illustrate a general point: when constraints are represented as rules, they can be queried, tested, audited, and monitored at run time. For computational-law treatments in a Prolog setting, see, for example, the discussion in \textit{Prolog: The Next 50 Years} and related work on executable norms and obligations \cite{PNFY2025,Kowalski2025}.

\subsubsection{Bridging, not siloing}

Hybrid architectures are not an either/or choice between Prolog and Python. Modern Prolog systems increasingly support tight, bidirectional interoperation with Python, and Janus in particular provides a fast two-way interface between Prolog and Python that is being developed as a de facto standard across multiple Prolog implementations. From an agent-design perspective, a common and often natural direction is for Prolog to call Python: the symbolic layer retains control and delegates vectorised numerics, deep-learning inference, or GPU-backed tensor algebra to Python libraries, while Prolog remains responsible for search, constraint enforcement, supervision, and auditability. This yields a pragmatic division of labour: Python for tensors, Prolog for control \cite{WielemakerJanus,SwiftAndersen2023}.

\subsection{Interaction patterns between LLMs and Prolog}
\label{sec:interaction-patterns-llm-prolog}

This subsection describes three progressively stronger integration styles. They are not mutually exclusive, and real systems often combine all three.

\subsubsection{LLM $\rightarrow$ Prolog (tool invocation)}

In a tool-augmented architecture, an LLM is embedded in an orchestrator that can route structured requests to external components. If we expose a Prolog engine as a tool, the LLM can decide when to delegate a subproblem to symbolic evaluation, and the result becomes a grounded intermediate value rather than a fabricated continuation. The essential design point is that the Prolog side must be sandboxed: allowed predicates are whitelisted, resource limits are enforced, and the result is returned in a typed, machine-checkable format.

Conceptually, one may imagine a helper predicate such as \texttt{prolog\_tool/2} that evaluates a restricted query and returns a structured answer term that the LLM can incorporate into its next step. This is the lightest-weight integration, and it already delivers substantial benefit when the delegated subtask has crisp semantics.

\subsubsection{Prolog $\rightarrow$ LLM (API delegation)}

Conversely, Web Prolog can delegate open-ended language tasks to an LLM and treat the model call as just another effectful predicate, guarded by timeouts, budgets, and policy checks. The exact API surface will evolve, but the architectural point remains stable: the symbolic layer controls when to call the model, how to constrain the prompt, how to validate the reply, and what to do when the call fails or returns something unsafe.

The following sketch shows the shape of an HTTP call to a contemporary text-generation endpoint. The details should be read as schematic and should be aligned with the provider's current documentation. Note that the option access is expressed using \texttt{options/3} rather than field access syntax.

\begin{verbatim}
:- use_module(library(http/http_client)).
:- use_module(library(http/http_json)).

call_llm(SystemMsg, UserMsg, Options, Reply) :-
    options(Options, endpoint, Endpoint),
    options(Options, authorization, Auth),
    options(Options, model, Model),
    Payload = json([
        model=Model,
        input=[
            json([role=system, content=SystemMsg]),
            json([role=user,   content=UserMsg])
        ]
    ]),
    http_post(Endpoint, json(Payload), Reply,
              [ request_header('Authorization'=Auth)
              ]).
\end{verbatim}

\noindent In a Web Prolog setting, this call is typically wrapped in a control shell: a timeout guard, a retry policy, a trace hook that records prompts and replies with provenance tags, and a validator that checks the reply against a schema before it re-enters the agent's control flow.

\subsubsection{Tight coupling and differentiable logic}

In tight neuro-symbolic systems, parts of the symbolic procedure participate directly in gradient-based training. The core idea is to replace crisp truth values with differentiable scores, and to compile proof search into a computation graph that can be optimized. Projects such as DeepProbLog and NeurASP demonstrate that this can be engineered in practice by compiling logic programs into weighted proof graphs or tensor expressions executed in a deep-learning framework \cite{DeepProbLog,NeurASP}.

From the Trinity perspective, the key observation is not that every agent should be differentiable, but that Prolog already provides a disciplined representation of structure, and that structure can sometimes be made learnable where it is genuinely useful, for example when learning soft constraints, priors, or perception modules. A Web Prolog node can still keep behavioural authority in explicit rules and policies, while allowing selected components to be optimized statistically.

\subsection{Statechart-wrapped medical triage agent}
\label{sec:statechart-wrapped-triage}

To make the cage pattern concrete, consider a simplified medical triage agent implemented as a statechart actor supervising an LLM.

\begin{enumerate}
\item A Web Prolog actor initialises a statechart with states \texttt{collect}, \texttt{confirm}, \texttt{diagnose}, \texttt{advise}, and \texttt{abort}.
\item In \texttt{collect}, the LLM generates patient-facing questions; answers are parsed by DCGs and stored as structured facts.
\item The transition \texttt{collect} $\rightarrow$ \texttt{confirm} fires when required slots such as \texttt{age/1} and \texttt{symptom/1} are instantiated.
\item In \texttt{diagnose}, Prolog rules map symptom sets to candidate conditions and compute risk scores under explicit assumptions.
\item If risk $\ge \theta$, the statechart transitions to \texttt{advise(escalate)}; otherwise to \texttt{advise(self\_care)}.
\item All model outputs pass through a validator before they re-enter the state machine, rejecting disallowed content (for example, prohibited advice classes, missing disclaimers, or policy violations) and forcing the agent into \texttt{abort} or \texttt{escalate} when appropriate.
\end{enumerate}

\noindent The central benefit is architectural: the LLM provides language, but the statechart provides behaviour. The system therefore has states, transitions, stopping conditions, and escalation paths, which are difficult to guarantee in a free-form prompt loop.

\subsection{Challenges and research directions}
\label{sec:neuro-symbolic-challenges}

\begin{itemize}
\item \textbf{Latency and budgeting.} When LLM calls sit on a critical interaction path, the agent must treat time and cost as first-class resources, using caching, local models, or co-location where appropriate, and degrading gracefully when budgets are exceeded.
\item \textbf{Tracing and audit.} Bridging symbolic traces (choice points, proof trees, constraint stores) with neural traces (prompts, tool calls, sampling settings) into a unified audit log remains a practical challenge, but it is also a major opportunity for web-native agent engineering.
\item \textbf{Proof artefacts as user-facing explanations.} Proof trees and derivation traces are precise but not always pleasant to read. An LLM could be used as a commentary layer that turns such traces into structured explanations (step-by-step narratives, minimal justifications, or counterfactual ``what would have to change'' analyses), while the Prolog artefact remains the authoritative source of truth.
\item \textbf{Intermediate representations.} Standardising typed schemas for the boundary between neural and symbolic modules reduces brittle prompt engineering. JSON with explicit types is often sufficient, but richer formats such as JSON-LD can help when provenance and linking matter.
\item \textbf{Safety envelopes.} A reliable agent needs a policy layer that is external to the model and enforceable by construction. In Web Prolog this naturally lives in the symbolic backplane and in the supervisory wrapper.
\item \textbf{Differentiability where it matters.} End-to-end gradient flow through parts of a logic program remains an active research direction, and it is most compelling when used selectively rather than ideologically.
\item \textbf{Traffic optimisation on the Prolog Web.} Because Web Prolog makes messaging and distribution explicit, it naturally exposes telemetry (latency, queue depths, error rates) that could support learning-based optimisation of higher-level traffic patterns on the Prolog Web, such as pacing, batching, replica selection, or load-aware routing. The architectural twist is that this can be done with hard safety and policy constraints enforced by the symbolic and actor-level substrate: learning chooses among allowed actions, while Prolog defines what is allowed \cite{Jay2019Deep,Xu2023Teal}.
\end{itemize}

\subsection{Summary: Toward a web-native neuro-symbolic MAS}
\label{sec:summary-toward-web-native-neuro-symbolic}

By placing Web Prolog at the heart of the agent architecture, we inherit constraints, symbolic search, actor-oriented concurrency, and supervision from the outset. LLMs then become pluggable neural co-processors rather than monolithic sources of truth. Compared with Python-centric stacks that treat control as an emergent property of prompts and glue code, the Web Prolog approach makes control a program. This is the core architectural claim of this section: neuro-symbolic agents are not merely LLMs with add-ons, but systems whose behaviour is defined by a symbolic substrate that can call neural components when language or perception is needed.

Hybrid architectures still need to communicate their state, assumptions, and constraints to humans. The most common place where these choices become user-visible is dialogue: turn-taking, repair, refusals, and escalation are precisely where a symbolic control envelope pays off. We therefore turn next to user-interface agents.

\section{User-interface agents}
\label{sec:embodied-agents}

\begin{quotation}
\flushright Content is king. -- \textit{Bill Gates} \\
\flushright Conversation is king. Content is just something to talk about. -- \textit{Cory Doctorow}
\end{quotation}

\vspace{0.5cm}

\noindent The Agentic Prolog Web is a substrate for multi-agent interaction in general, but many of the most common and economically important agents are not peers in a negotiation protocol. They are user-interface agents: components that mediate between a person and a computational environment. In such settings, the engineering problem is not only what the agent can infer, but how it manages an interaction over time: when to ask, when to confirm, how to repair misunderstandings, how to handle silence and timeouts, and when to refuse or escalate.

This section therefore shifts emphasis from agent-to-agent coordination to dialogue as control. The central claim is that conversational behaviour should not be treated as an emergent property of a prompt. It should be treated as a program with states, transitions, and failure handling. Web Prolog is well suited to this stance because it can combine (i) symbolic representations of dialogue-relevant facts, (ii) executable parsing and interpretation, and (iii) an actor-oriented control shell, most naturally expressed as a statechart.

\subsection{UI agents versus MAS agents}
\label{sec:ui-agents-versus-mas}

In classic multi-agent systems, interaction is often modelled as peer-to-peer exchange: negotiation, coordination, and distributed problem solving among relatively symmetric participants. User-interface agents are asymmetric. They exist to serve a person, and they are evaluated by human-facing criteria such as clarity, trust, timing, politeness, and error recovery. The success condition is not only that the agent eventually produces a correct plan, but that it maintains an intelligible and safe interaction while doing so.

On the Web, this role frequently manifests as a conversational web agent: an agent that accepts input in natural language (text or speech) and produces responses that guide the user through a task. In such systems, the boundary between reasoning and interaction is not clean. Questions, confirmations, clarifications, and refusals are themselves action choices that must be governed by policy.

\subsection{Turn-taking and repair as control structure}
\label{sec:turn-taking-and-repair}

Conversation, in the technical sense used in linguistics, is a jointly managed activity. Participants alternate turns under a largely implicit turn-taking system; many contributions occur in adjacency-pair patterns (question--answer, offer--acceptance, request--grant); and when misunderstanding occurs, participants engage in repair sequences to re-establish common ground.

For an engineered UI agent, the practical implication is that interaction cannot be reduced to a single request--response call. The agent needs persistent state and control over what may happen next. It must be able to suspend a task to ask a question, resume after an answer, back up after a misunderstanding, and time out or escalate when interaction fails. These are precisely the kinds of behaviours that state machines and statecharts are designed to express.

A useful mental model is to separate the conversational system into a \emph{flow layer} and a \emph{content layer}. The flow layer governs when the system may advance, when it must confirm, what to do on silence, and how to recover after a misunderstanding. The content layer is where parsing, belief update, and reasoning live. In a Web Prolog setting, both layers can be expressed as programs, but the flow layer benefits from an explicit statechart representation because the control surface is the point where reliability and user trust are won or lost.

\subsection{Dialogue as proof: the missing-axioms view}
\label{sec:missing-axioms-theory}

A useful way to connect dialogue management to symbolic AI is the Missing Axioms Theory (Smith \& Hipp, 1994). In this view, the agent's task can be phrased as an attempt to prove a goal. When the proof fails, the reason is often not that the rules are wrong, but that some required facts are missing. Dialogue then becomes a controlled process for acquiring the missing axioms.

This perspective is already latent in the askable expert-system pattern from Section~\ref{sec:query-the-user} (and, in web-native form, Section~\ref{sec:expert-system-agent}). There, proof search determines which question must be asked next. The limitation of the toy meta-interpreter approach is not the idea, but the control envelope: real dialogue involves interruptions, repairs, topic shifts, and timeouts. To handle those robustly, we need an interaction controller with persistent state and well-defined failure handling.

\subsection{Statecharts as dialogue governors}
\label{sec:statecharts-as-dialogue-governors}

A practical architecture is to separate flow from content. The flow concerns turn-taking, timeouts, confirmations, and repair. The content concerns parsing, storing extracted facts, and reasoning about constraints and next steps. Statecharts are a good fit for the flow layer because they support hierarchy and orthogonality, events, guarded transitions, actions, and a history mechanism.

In Web Prolog, a statechart actor can therefore serve as the authoritative dialogue governor, while Prolog code in the statechart's data model provides the symbolic substrate for interpretation and reasoning. A minimal illustration is an agent that collects two attributes (colour and size) and uses a history state to recover after a misunderstanding. Operationally, this means that the dialogue manager needs a representation of where it is in the interaction, what information is still missing, and what repair strategies are permitted.

\begin{verbatim}
<statechart datamodel="web-prolog" initial="Start">
  <datamodel>
    :- dynamic color/1, size/1.
    np(Attr) --> [a], adj(Attr), [apple].
    adj(color(red)) --> [red].
    adj(size(big))  --> [big].
  </datamodel>

  <state id="Start" initial="Color">
    <history id="H">
      <go to="Color"/>
    </history>
    <state id="Color">
      <onentry>
        writeln('What color do you want?')
      </onentry>
      <go to="Size" 
          on="input(Input)" 
          if="phrase(np(color(C)), Input)">
        assert(color(C))
      </go>
    </state>
    <state id="Size">
      <onentry>
        writeln('What size do you want?')
      </onentry>
      <go to="Final" 
          on="input(Input)" 
          if="phrase(np(size(S)), Input)">
        assert(size(S))
      </go>
    </state>
    <go to="H" on="input(_)">
      writeln('I did not understand.')
    </go>
  </state>
  <final id="Final">
    <onentry>
      size(S), color(C),
      format('You want a ~p ~p apple!~n', [S,C])
    </onentry>
  </final>
</statechart>
\end{verbatim}

\noindent Three elements matter here.

\subsubsection{Repair via history}

The history state supports a simple repair policy for unrecognised input. When an \texttt{input(\_)} event fails to match any more specific transition, the machine emits a complaint and transitions to the history pseudostate. This re-enters the most recently active substate of the dialogue, causing its \texttt{<onentry>} prompt to run again and thereby re-asking the current question. In contrast to restarting the whole interaction, the repair stays within the current plan, resembling how human repair sequences typically preserve the larger conversational trajectory.

\subsubsection{Parsing as a guard}

Transition guards call DCG parsing directly, so the controller advances only when the incoming \texttt{input(Input)} can be parsed into the intended structured form. This makes the interface between interpretation and control explicit: the grammar determines what counts as understood, while the statechart determines when understanding is sufficient to commit to a control move.

\subsubsection{Belief updates as explicit effects}

Executable content makes belief updates explicit and localised. When an input is successfully parsed, the extracted fact is asserted as a named predicate in the statechart's data model. These stored facts can later be queried, audited, and constrained by ordinary Prolog rules, rather than remaining implicit in conversational context.

\medskip
\noindent Together, history-based repair, guard-based interpretation, and explicit data-model updates yield a compact pattern for dialogue control: recognise \(\rightarrow\) commit \(\rightarrow\) store, otherwise complain \(\rightarrow\) repeat.

\subsection{Handover: dialogue as a governed interface}
\label{sec:ui-agents-handover}

Statechart-governed dialogue provides predictability, repair structure, stopping conditions, and a principled place to implement refusals and escalation. In the neuro-symbolic setting developed in Section~\ref{sec:neuro-symbolic}, this control envelope plays an additional role: it defines the boundary between what the language model may suggest and what the agent may do.

Seen this way, the statechart is not merely a dialogue manager. It is a governor for an LLM-augmented interface: the model can be invoked for paraphrasing, interpretation, and question generation, but its outputs re-enter the system only through guards, schemas, and policies that are authoritative and auditable. Dialogue thus becomes the most visible instance of the general design stance of this chapter: place behavioural authority in the symbolic substrate, and use neural components as pluggable co-processors where language robustness is needed.

\section{Summary}
\label{sec:summary-and-discussion}

This chapter framed Prolog's relevance in an AI-driven software landscape as an operational question: can we make symbolic constraints, search, and explanations deployable and composable at Internet scale, so that they can supervise more fallible components in agentic systems. We first recalled Prolog's symbolic competence on a single machine: representing knowledge as executable clauses, exploring discrete spaces systematically, planning as search, constructing proof objects for explanation, acquiring missing axioms through interaction, and using DCGs to connect language and structure. We then lifted these idioms onto the Agentic Web, treating Prolog engines as networked actors and letting proofs and explanations span nodes.

On that foundation, we argued for a neuro-symbolic stance: treat large language models as powerful but fallible language modules, embed them in a symbolic control and constraint envelope, and make behavioural authority explicit. We noted that the ecosystem changes rapidly, and that even code generation itself may become increasingly automated, but that this makes precise semantics and explicit invariants more valuable rather than less.

Finally, we turned to user-interface agents as the most common and user-visible application class of these ideas. Dialogue is where control matters: turn-taking, repair, confirmation, timeouts, refusals, and escalation must be governed by something more concrete than a prompt. Statecharts provide this governance, and in Web Prolog they can be executed as actors whose guards and effects are symbolic programs. In that setting, the LLM becomes a co-processor for language, while the statechart and the Prolog substrate remain responsible for what the agent is allowed to do, what it knows, and why it believes what it believes.




